
\documentclass[a4paper,fleqn]{cas-sc}

\usepackage[authoryear,longnamesfirst]{natbib}
\usepackage{mathtools}
\usepackage{siunitx}
\usepackage{verbatim}

\newcommand{\gcm}{\,\mathrm{g\,cm^{-2}}}
\newcommand{\mGyday}{\,\mathrm{mGy\,day^{-1}}}
\newcommand{\mSvday}{\,\mathrm{mSv\,day^{-1}}}
\newcommand{\keVum}{\,\mathrm{keV\,\mu m^{-1}}}
\newcommand{\Qeff}{Q_\mathrm{eff}}
\newcommand{\LETd}{\overline{L}_D}
\newcommand{\REID}{\mathrm{REID}}
\newcommand{\code}[1]{\texttt{#1}}

\begin{document}
\let\WriteBookmarks\relax
\def\floatpagepagefraction{1}
\def\textpagefraction{.001}

\shorttitle{Supplementary Material: Open-source GCR dosimetry pipeline}
\shortauthors{A.H.~Shah}

\title[mode=title]{Supplementary Material for: An open-source GCR
dosimetry pipeline for Mars mission risk assessment: organ-specific REID,
SEP acute hazard, and global uncertainty quantification}

\author[1]{Aryan Hemendra Shah}[orcid=0009-0008-0518-428X]
\cormark[1]
\ead{ahs222@miami.edu}

\affiliation[1]{organization={College of Arts and Sciences},
            addressline={University of Miami},
            city={Miami},
            state={FL},
            postcode={33146},
            country={USA}}

\cortext[1]{Corresponding author}

\begin{abstract}
This document contains full physics/biology derivations, extended parameter
tables, and extended discussion supporting the main text, which is
constrained to 3500 words (Introduction through Discussion) per the
\emph{Life Sciences in Space Research} author guidelines. Section, table,
and equation numbers here are independent of the main text and prefixed
with S.
\end{abstract}

\maketitle

\tableofcontents

\section{Supplementary Methods}

\subsection{Mission Trajectory Model}

A Hohmann-transfer trajectory from Earth (1.0~AU) to Mars (1.524~AU) is
generated by computing heliocentric distance as a function of mission elapsed
time, $r(t)$, using Keplerian orbital mechanics:
\begin{equation}
  r(t) = \frac{a(1-e^2)}{1 + e\cos\theta(t)},
\end{equation}
where $a = (r_\oplus + r_\mathrm{Mars})/2 = 1.262$~AU is the semi-major axis
and $e = (r_\mathrm{Mars} - r_\oplus)/(r_\mathrm{Mars} + r_\oplus) = 0.208$
is the eccentricity of the transfer ellipse. Transit duration is fixed at
259~days (one half-period of the transfer orbit). The heliocentric distance
enters the dose rate only through the solar modulation potential $\phi(t)$;
direct geometric flux dilution as $r^{-2}$ is not applied separately because
the force-field modulation already encodes the energy dependence of solar
wind effects at the relevant energies.

\subsection{Solar Modulation Potential}

The time-resolved solar modulation potential $\phi(t)$ is obtained from the
Usoskin et~al.\ (2017) database, which provides monthly $\phi$ values
reconstructed from ground-level neutron monitor count rates back to 1936.
For the MSL transit period (2011 November -- 2012 July), values cluster
around $\phi \approx 481 \pm 50$~MV. Systematic offsets between different
neutron monitor stations introduce $\pm 50$--$100$~MV uncertainty in the
absolute $\phi$ calibration; this is propagated via the \code{phi\_scale}
Latin Hypercube Sampling (LHS) parameter (Table~\ref{tab:uq_params}).

\subsection{GCR Spectrum Generation}

\subsubsection{Force-Field Modulation}

Solar modulation of the galactic cosmic ray (GCR) flux is computed using
the Gleeson--Axford (1968)
force-field approximation. For an ion of charge $Z$ and mass number $A$, the
modulated differential flux at kinetic energy per nucleon $E$ is:
\begin{equation}
  J(E, \phi) = J_\mathrm{LIS}(E + \delta E)
  \cdot \frac{E(E + 2E_0)}{(E + \delta E)(E + \delta E + 2E_0)},
  \label{eq:ff}
\end{equation}
where $E_0 = 938.272$~MeV is the proton rest-mass energy,
$\delta E = (Z/A)\,\phi$ (MeV/nucleon), and $J_\mathrm{LIS}$ is the local
interstellar spectrum (LIS, below). The ratio in Eq.~\eqref{eq:ff} is the
Jacobian of the energy transformation, ensuring flux conservation. The
approximation is accurate for rigidities above $\sim\!1$~GV and $\phi$ in
$200$--$1500$~MV \citep{potgieter2013}; charge-sign-dependent drift (relevant
at solar minimum) is not captured, but dose here is dominated by particles
above a few hundred MeV/n where this is not expected to matter materially.

\subsubsection{Local Interstellar Spectra}

Proton and helium LIS follow Vos \& Potgieter \cite{vos2015}:
\begin{equation}
  J_\mathrm{LIS}^\mathrm{H}(E) = C_\mathrm{H}
  \cdot \frac{(E_\mathrm{GeV} + 0.67)^{-3.93}}{\beta^2},
\end{equation}
with $C_\mathrm{H} = 1.9 \times 10^4$~cm$^{-2}$~s$^{-1}$~MeV$^{-1}$~sr$^{-1}$.
Heavy ions (C, O, Si, Fe) follow the per-species HelMod model of Boschini
et al.\ \cite{boschini2020}:
\begin{equation}
  J_\mathrm{LIS}^{(s)}(E_n) = C_s \cdot
  \frac{(E_\mathrm{GeV/n} + \delta_s)^{-\alpha_s}}{\beta^2}.
\end{equation}
Table~\ref{tab:lis_params} lists all adopted parameters. Iron's flatter
index ($\alpha_\mathrm{Fe}=2.65$) gives it disproportionate high-energy flux,
amplified by $Z^2$ stopping-power scaling.

\begin{table}[t]
  \centering
  \caption{LIS spectral parameters used in \code{gcrisk/spectrum.py}.
           Proton and helium from Vos \& Potgieter \cite{vos2015}; heavy ions from
           Boschini et al.\ \cite{boschini2020}.}
  \label{tab:lis_params}
  \small
  \begin{tabular*}{\tblwidth}{@{}LCCCL@{}}
    \toprule
    Species & $C_s$ & $\alpha_s$ & $\delta_s$ (GeV/n) & Primary source \\
    \midrule
    H       & $1.90\times10^4$ & 3.93 & 0.67 & Vos \& Potgieter (2015) \\
    He      & $1.82\times10^3$ & 2.78 & 0.62 & Vos \& Potgieter (2015) \\
    C       & $2.50\times10^2$ & 2.72 & 0.50 & Boschini et al.\ (2020) \\
    O       & $1.80\times10^2$ & 2.69 & 0.50 & Boschini et al.\ (2020) \\
    Si      & $3.00\times10^1$ & 2.70 & 0.50 & Boschini et al.\ (2020) \\
    Fe      & $4.30\times10^0$ & 2.65 & 0.50 & Boschini et al.\ (2020) \\
    \bottomrule
  \end{tabular*}
\end{table}

\subsubsection{Composition-Constrained Normalization}
\label{sec:calib_supp}

Calibration is a two-step procedure. \textbf{Step 1 (composition
constraint)}: for each species $s$, a correction factor
$r_s = (J_s/J_\mathrm{H})_\mathrm{ACE/CRIS} / (J_s/J_\mathrm{H})_\mathrm{LIS}$
is computed so the modulated LIS reproduces the particle-flux ratio measured
by ACE/CRIS \citep{george2009} and PAMELA \citep{adriani2014} at
$E_n=200$~MeV/n, $\phi=481$~MV (Table~\ref{tab:ace_cris_comp}).
\textbf{Step 2 (global dose scale)}: a single global scale
$k = 0.894$ is then fit so $J_s \leftarrow k \cdot r_s \cdot J_s$ reproduces
$\dot{D}=1.84\mGyday$ at $16\gcm$ Al, $\phi=481$~MV \citep{zeitlin2013}. This
replaces a prior six-free-parameter fit (one factor per species fit
independently to the single dose scalar, with no physical composition
constraint) with one free parameter whose relative species mix is
externally anchored.

\begin{table}[t]
  \centering
  \small
  \caption{ACE/CRIS and PAMELA reference particle-flux ratios relative to
           $J_\mathrm{H}=1.0$, $E_n \approx 200$~MeV/n, $\phi \approx 481$~MV
           (2011--2012 MSL epoch).}
  \label{tab:ace_cris_comp}
  \begin{tabular*}{\tblwidth}{@{}LCL@{}}
    \toprule
    Species & $J_s / J_\mathrm{H}$ & Source \\
    \midrule
    He & $6.60\times10^{-2}$ & Adriani et al.\ (2014) \\
    C  & $1.80\times10^{-3}$ & George et al.\ (2009) \\
    O  & $2.20\times10^{-3}$ & George et al.\ (2009) \\
    Si & $1.80\times10^{-4}$ & George et al.\ (2009) \\
    Fe & $1.40\times10^{-4}$ & George et al.\ (2009) \\
    \bottomrule
  \end{tabular*}
\end{table}

Table~\ref{tab:species_compare} compares the resulting composition against
the NASA Space Radiation Laboratory (NSRL) GCR simulator reference field
\citep{slaba2016}. The proton fraction is within $\sim\!4$ points of NSRL;
the high-charge, high-energy (HZE) fraction is $\sim\!1.7\times$ the NSRL
reference (down from
$\sim\!4\times$ under the prior calibration), attributed to
continuous-slowing-down approximation (CSDA) transport
not modeling nuclear fragmentation (which converts high linear-energy-transfer
(LET) heavy ions into
lower-LET secondaries in NSRL/HZETRN but not here). The helium fraction is
below NSRL, reflecting the measured ACE/CRIS He/H ratio at $200$~MeV/n being
lower than the LIS model's He/H at the higher energies that set the
unmodulated normalization.

\begin{table}[t]
  \centering
  \small
  \caption{Per-charge-group absorbed dose fractions: composition-constrained
           pipeline vs.\ NSRL reference \citep{slaba2016}. Geometries differ
           by $4\gcm$ Al and solar phase; evaluating the pipeline at NSRL
           conditions ($20\gcm$ Al, $\phi\approx300$~MV) gives
           $\mathrm{HZE}=12.6\%$ vs.\ NSRL's $7.6\%$.}
  \label{tab:species_compare}
  \begin{tabular*}{\tblwidth}{@{}LCC@{}}
    \toprule
    Particle class & Pipeline ($16\gcm$ Al, $\phi=481$) & NSRL ($20\gcm$ Al, solar min) \\
    \midrule
    $Z=1$ (protons)  & 68.9\% & 73.3\% \\
    $Z=2$ (helium)   &  9.8\% & 19.2\% \\
    $Z>2$ (HZE)      & 13.0\% &  7.6\% \\
    \bottomrule
  \end{tabular*}
\end{table}

With this composition, the pipeline gives $\dot{H}=6.09\mSvday$ and
$\Qeff=3.32$ at the calibration geometry, vs.\ RAD's $4.81\pm0.97\mSvday$
and $2.62\pm0.14$ (a $27\%$ excess). This is attributed to the same CSDA
fragmentation gap: without converting high-LET HZE dose into lower-LET
secondaries, the dose-weighted quality factor is inflated. $\dot{H}$ and
organ-specific cancer risk (REID) from this pipeline should be read as
conservative upper estimates.
The \code{hze\_norm\_scale} LHS parameter (LogNormal, $\sigma=0.20$)
propagates the residual $\sim\!20\%$ ACE/CRIS inter-instrument composition
uncertainty \citep{george2009, boschini2020}.

\subsection{Charged Particle Transport}

\subsubsection{Stopping Power and CSDA Range}

Energy loss is computed via CSDA. Stopping power for charge $Z$, velocity
$\beta$:
\begin{equation}
  S(E/A,\, Z) = z_\mathrm{eff}^2(E/A,\, Z) \cdot S_p(E/A),
\end{equation}
with Barkas effective charge \citep{barkas1963}
$z_\mathrm{eff} = Z[1-\exp(-125\beta Z^{-2/3})]$ (reducing to $Z^2$ scaling
above $\sim\!1$~GeV/n). Proton stopping powers (Al, PE, water, tissue) are
tabulated from NIST PSTAR/ASTAR (500 log-spaced points, 1~MeV--$10^5$~MeV,
cubic-spline interpolated). CSDA range $R(E)=\int_0^E dE'/S(E',Z)$; exit
energy after slab $x$ is $E_\mathrm{out}=R^{-1}(R(E_\mathrm{in})-x_\mathrm{eff})$
with $x_\mathrm{eff}=x\cdot Z^2/A$; particles with $R(E_\mathrm{in})<x_\mathrm{eff}$ stop.

\subsubsection{Nuclear Interaction Attenuation}

Inelastic nuclear interactions are modeled as exponential fluence
attenuation, $\Phi_s(E,x)=\Phi_s^{(0)}(E)\exp(-x/\lambda_s)$, with mean
free path $\lambda_s$ from the Bradt-Peters geometric cross-section
\citep{bradtpeters1950}, $\sigma=\pi r_0^2(A_p^{1/3}+A_t^{1/3}-\delta)^2$
($r_0=1.35$~fm, $\delta=0.83$). This is the same geometric form used to build
HZETRN's tabulated mean free paths \citep{wilson1991}. This does not model
the secondary fragment spectrum; above $\sim\!30\gcm$ for Fe/Si this
approximation becomes increasingly crude as fragmentation produces
biologically-relevant lighter secondaries (C, O). A full nuclear reaction
network (HZETRN, Geant4/TOPAS) would be needed to close this gap. This is the
primary physical limitation of this transport model.

\subsubsection{Secondary Neutron Production}

Rather than tracking individual neutron histories, neutron dose equivalent
uses a pre-computed lookup, $\dot{H}_n(x,\phi,\mathrm{mat}) =
\dot{H}_n^\mathrm{Al}(\phi)\cdot f_\mathrm{mat}\cdot g(x)$, with material
yield factor $f_\mathrm{mat}$ (PE $\approx 0.72\times$ Al) and depth
dependence $g(x)$ interpolated from the HZETRN benchmark of Slaba et al.\
\cite{slaba2014}; neutron $Q_n=10$ (International Commission on
Radiological Protection Publication 60, ICRP-60). This introduces
$\pm25$--$30\%$ uncertainty, propagated via the \code{neutron\_H} LHS
parameter.

\subsection{Organ Dose Model}

Organ dose is computed by transporting shielded flux through an additional
tissue layer at each organ's ICRP Publication~89 depth \citep{icrp89}
(Table~\ref{tab:organ_depths}); total areal density at organ $k$ is
$x_\mathrm{total}^{(k)} = x_\mathrm{wall}+d_k$. Absorbed dose rate:
\begin{equation}
  \dot{D}_k = \sum_s \int \Phi_s(E; x_\mathrm{total}^{(k)},\phi)\cdot S_s(E)\cdot\rho^{-1}\cdot C\,dE,
\end{equation}
with the $4\pi$ isotropic-flux factor absorbed into $C=1.602\times10^{-10}$~Gy
per (MeV g$^{-1}$cm$^{-2}$); energy grid is 200 log-spaced points,
$10$--$10^5$~MeV/n ($<0.5\%$ convergence).

\begin{table}[t]
  \centering
  \caption{Organ tissue depths and ICRP-60 tissue weighting factors.
           Depths from ICRP Publication~89 \citep{icrp89}.}
  \label{tab:organ_depths}
  \small
  \begin{tabular*}{\tblwidth}{@{}LCCL@{}}
    \toprule
    Organ & Depth (g/cm$^2$) & $w_T$ (ICRP-60) & Cancer type modeled \\
    \midrule
    Stomach    & 8.0  & 0.12 & Gastric carcinoma \\
    Colon      & 9.5  & 0.12 & Colorectal carcinoma \\
    Lung       & 5.0  & 0.12 & Lung carcinoma \\
    Leukemia   & 4.0  & 0.12 & Myeloid leukemia \\
    Esophagus  & 4.0  & 0.05 & Esophageal carcinoma \\
    Liver      & 6.0  & 0.04 & Hepatocellular carcinoma \\
    Bladder    & 10.0 & 0.04 & Transitional cell carcinoma \\
    Thyroid    & 2.0  & 0.04 & Thyroid carcinoma \\
    Breast     & 1.5  & 0.05 & Breast adenocarcinoma \\
    Ovary      & 12.0 & 0.20 & Ovarian carcinoma \\
    Other      & 7.0  & 0.05 & Composite remainder \\
    \bottomrule
  \end{tabular*}
\end{table}

Dose equivalent: $\dot{H}_k = \sum_s Q_s \cdot \dot{D}_{k,s}$, with ICRP-60
$Q(L)$ \citep{icrp60}:
\begin{equation}
  Q(L) = \begin{cases}
    1                & L < 10\keVum \\
    0.32\,L - 2.2   & 10 \leq L \leq 100\keVum \\
    300 / \sqrt{L}  & L > 100\keVum
  \end{cases}
  \label{eq:quality}
\end{equation}
applied per species at the dose-averaged LET, $Q_s=Q(\overline{L}_{D,s})$.
This mean-LET approximation (vs.\ integrating $Q(L)$ over each species' full
LET spectrum) underestimates $Q$ for carbon by $\sim\!16\%$ and overestimates
it for oxygen by $\sim\!24\%$ (both straddle the $10\keVum$ kink), with
$\lesssim3\%$ bias for H, He, Si, Fe; the opposite-signed C/O biases
partially cancel in the field-averaged $\Qeff = \sum H_{k,s}/\sum D_{k,s}$,
which is $2.60$ at the calibration geometry (before the neutron
contribution raises the total apparent $\Qeff$ to $3.32$), within $0.8\%$
of the RAD ion-only estimate. This is expected, since it follows directly
from the calibration procedure.

\subsection{Cancer Risk Model (REID) — Full Formulation}

REID integral:
\begin{equation}
  \mathrm{REID}(a_e, s) = \int_{a_e + L}^{a_\mathrm{LT}}
    \frac{\mu_c(a, s) \cdot S(a \mid a_e, s)}{S(a_e \mid s)} \, da,
  \label{eq:reidintegral}
\end{equation}
$a_e$=age at exposure, $L=10$~yr minimum latency, $S(a|a_e,s)$=conditional
survival. Excess mortality combines the excess absolute risk (EAR) and
excess relative risk (ERR):
\begin{equation}
  \mu_c(a) = \lambda_c^\mathrm{background}(a) \cdot \mathrm{ERR}(D, a_e)
           + \mathrm{EAR}(D, a_e, a),
\end{equation}
with EAR/ERR organ-specific functions of dose, age at exposure, and attained
age from Tables 2--3 of Cucinotta et al.\ \cite{cucinotta2013}. The dose and
dose-rate effectiveness factor (DDREF) reduces
low-dose-rate risk, $D_\mathrm{eff}=D/\mathrm{DDREF}$: point estimates use a
fixed DDREF of $1.75$; the uncertainty ensemble instead samples
DDREF~$\sim$~Uniform$[1,2]$. Conditional survival and baseline cancer
mortality use interpolated US sex-specific lookup tables in
\code{gcrisk/reid.py}. This is a practical approximation to the NASA/Cucinotta
framework.

The REID distribution is right-skewed, driven by \code{Q\_factor}'s
heavy-tailed lognormal draws: its LogNormal[1.0,$\ln 2$] draws reach a
sampled maximum of $\sim\!8.2\times$
in the $N=500$ ensemble, corresponding to $\Qeff\approx27$ at the
calibration point. This approaches but does not exceed the ICRP-60 physical
maximum of $30$ ($L=100\keVum$). Every LHS sample's \code{Q\_factor} is
therefore clipped so effective $Q$ cannot exceed $30$; for the reported
$N=500$ ensemble this clip does not bind, so reported percentiles are
unchanged, but it removes any possibility of an unphysical excursion at
larger $N$ or different seeds.

REID here is computed as a direct hazard-rate integral
(Eq.~\ref{eq:reidintegral}). A survival-curve formulation like
$1-\exp(-\mu)$ would confine individual Monte Carlo samples to $[0,1]$
automatically; this direct-integral form does not. Each sample is
therefore also clipped to $[0,1]$ post hoc, since REID represents a
probability and cannot exceed $100\%$ by definition. This second clip
barely binds at the reported calibration-geometry ensemble (1 of 500
samples, unclipped maximum $101.3\%$) and does not move any reported
percentile, but it removes draws that are more frequently, and more
severely, unphysical at low shielding (REID vs.\ shielding figure, main
text), where the unshielded dose is large enough that the heavy-tailed
\code{Q\_factor} draws alone could otherwise push individual samples to
$200$--$300\%$.

\subsection{Solar Energetic Particle Acute Risk}

Solar energetic particle (SEP) fluence spectra use the Band function
\citep{tylka2006}:
\begin{equation}
  J_\mathrm{SEP}(E) = C \times \begin{cases}
    E^{\gamma_\ell} \exp(-E/E_b) & E \leq E_\mathrm{trans}, \\
    \left[(\gamma_\ell - \gamma_h) E_b\right]^{\gamma_\ell - \gamma_h}
    \exp(\gamma_h - \gamma_\ell) \cdot E^{\gamma_h} & E > E_\mathrm{trans},
  \end{cases}
\end{equation}
$E_\mathrm{trans}=(\gamma_\ell-\gamma_h)E_b$. Three canonical events
(Table~\ref{tab:sep_events}) use parameters from Tylka \& Dietrich
\cite{tylka2006}; Aug~1972 predates space-borne dosimetry and is
reconstructed from ground-level neutron monitor/riometry data, carrying
larger uncertainty ($\sim\pm50\%$) than the 2003/2005 events.

\begin{table}[t]
  \centering
  \caption{Band-function spectral parameters for three canonical SEP events.}
  \label{tab:sep_events}
  \small
  \begin{tabular*}{\tblwidth}{@{}LCCCCC@{}}
    \toprule
    Event & $C$ & $\gamma_\ell$ & $\gamma_h$ & $E_b$ (MeV) & Duration (days) \\
    \midrule
    Aug 1972 & $4.0\times10^4$ & $-1.78$ & $-3.98$ & 32  & 2 \\
    Oct 2003 & $1.5\times10^4$ & $-1.40$ & $-3.20$ & 50  & 3 \\
    Jan 2005 & $3.0\times10^2$ & $-0.90$ & $-2.80$ & 110 & 1 \\
    \bottomrule
  \end{tabular*}
\end{table}

Blood-forming organ (BFO) absorbed dose transports event-integrated proton
fluence through wall +
$10\gcm$ tissue via CSDA; only protons are tracked (heavy-ion SEP fluence is
generally dose-negligible at these energies). Results compare to the NASA
30-day BFO limit of $250$~mGy \citep{cucinotta2010}. Mission encounter
probability is Poisson, $P_\mathrm{SEP}(T)=1-\exp(-\lambda T/365.25)$,
$\lambda=0.8$~events/yr at solar maximum, dropping $3$--$5\times$ at solar
minimum \citep{feynman1993}.

\subsection{Uncertainty Quantification via Latin Hypercube Sampling}

Nine parameters are varied via LHS (Table~\ref{tab:uq_params}):
five physics (solar modulation, LIS normalization, nuclear cross-sections,
neutron yield, HZE composition) and four biology (quality factor, DDREF,
ERR scale, EAR scale).

\begin{table}[t]
  \centering
  \caption{Latin Hypercube Sampling parameters, distributions, and physical
           justifications. LogNormal[$1.0$, $\sigma$] means a multiplicative
           factor $\exp(\mathcal{N}(0,\sigma))$ with median $1.0$.}
  \label{tab:uq_params}
  \small
  \begin{tabular*}{\tblwidth}{@{}LLL@{}}
    \toprule
    Parameter & Distribution & Justification \\
    \midrule
    \code{phi\_scale}      & LogNormal[1.0, 0.15] & Usoskin NM inter-station spread ($\sim\!15\%$) \\
    \code{LIS\_norm}       & LogNormal[1.0, 0.12] & Voyager/AMS-02 calibration spread ($\sim\!12\%$) \\
    \code{cross\_sec}      & LogNormal[1.0, 0.20] & Nuclear $\sigma$ tables: $\sim\!20\%$ \\
    \code{neutron\_H}      & LogNormal[1.0, 0.25] & HZETRN neutron table: $\sim\!25\%$ \citep{slaba2014} \\
    \code{hze\_norm\_scale}& LogNormal[1.0, 0.20] & ACE/CRIS inter-instrument spread \citep{george2009, boschini2020} \\
    \code{Q\_factor}       & LogNormal[1.0, $\ln 2$] & ICRP-60 $Q(L)$ form uncertainty: factor-of-2 \\
    \code{DDREF}           & Uniform[1.0, 2.0] & BEIR VII / ICRP-103 range \\
    \code{ERR\_scale}      & Normal[1.0, 0.30], clipped at 0 & Cucinotta (2013) ERR fit residuals \\
    \code{EAR\_scale}      & Normal[1.0, 0.35], clipped at 0 & Cucinotta (2013) EAR fit residuals \\
    \bottomrule
  \end{tabular*}
\end{table}

Lognormal parameters cannot drive their model input negative; \code{ERR\_scale}
and \code{EAR\_scale} are clipped at zero (affects $<0.1\%$ of draws, no
visible effect on percentiles). For each LHS sample, a full mission dose
integration is performed (spectrum, transport, organ dose, quality factor,
REID). Publication results use $N=500$ ($\sim\!90$ min, single core,
unparallelized); automated tests use $N=20$
(\code{UNCERTAINTY\_SAMPLES} env var).

Sensitivity is estimated via normalized Spearman $\rho^2$,
$S_i=\rho_i^2/\sum_j\rho_j^2$. This is an approximation to first-order Sobol
indices that underestimates interaction effects but is tractable at modest
sample counts. A proper Sobol analysis would need $\sim\!10^4$ samples
\citep{saltelli2010}.

The value-of-information sweep re-runs the $N=500$ ensemble once per
parameter, halving that parameter's spread (sigma, or DDREF's range width
centered on the same mean) while holding all others at baseline
(\code{narrow} argument to \code{gcrisk.uncertainty.lhs\_samples},
\url{scripts/sensitivity_narrowing_sweep.py}), recording the resulting
change in p95 REID.

\subsection{Dose-Averaged LET and Proton RBE Self-Consistency Check}

Dose-averaged LET per species: $\overline{L}_{D,s} = \int\Phi_sS_s^2\,dE
/ \int\Phi_sS_s\,dE$; field value $\overline{L}_D=\sum_sf_s\overline{L}_{D,s}$,
$f_s=D_s/\sum_{s'}D_{s'}$.

The repository also contains a proton relative biological effectiveness
(RBE) self-consistency layer using this
same $\overline{L}_D$ infrastructure: Wedenberg et~al.\ (2013)
\citep{wedenberg2013} and McNamara et~al.\ (2015) \citep{mcnamara2015},
matched to OpenTOPAS-RBE scorer conventions and checked against
OpenTOPAS-RBE reference values for a $100$~MeV proton beam in water (V79,
$\alpha/\beta=1.41$~Gy). This is an implementation self-consistency check,
and the foundation for a planned separate proton-therapy crossover
study.

\section{Supplementary Discussion: Full Limitations}

\subsection{Transport Model Approximations}
\begin{itemize}
  \item \textbf{Nuclear fragmentation secondaries.} Heavy-ion nuclear
        interactions are treated as simple fluence removal. As such, the pipeline doesn't take into account the lighter-ion secondaries produced through nuclear fragmentation.  This leads to the secondary
        proton/helium flux being underestimated at depth and overestimates iron attenuation,
        largest above $30\gcm$. Based on the general direction/magnitude of
        differences between simplified transport and full fragmentation
        codes \citep{wilson1991, slaba2014}, my order-of-magnitude estimate
        is that omitting secondary production overestimates $\Qeff$ at
        $16\gcm$ by roughly $15$--$25\%$. This is my own rough
        extrapolation from those sources' general direction and magnitude.
        The resulting REID overestimate is likely similar in fractional
        magnitude, well within the $12\times$ p5--p95 band but not
        negligible relative to the median.
  \item \textbf{Delta-ray production.} Secondary electrons from ion tracks
        are not tracked; small correction for macroscopic dose, would matter
        for microscopic track-structure DNA-damage models.
  \item \textbf{Geometry simplification.} Spacecraft modeled as a uniform
        aluminum slab; real geometries need ray-tracing (OLTARIS, Geant4).
  \item \textbf{Stopping power uncertainties.} NIST PSTAR/ASTAR accurate to
        $\sim\!1$--$2\%$ above 1 MeV/n, less certain at lower energies
        (shell corrections).
\end{itemize}

\subsection{Spectrum and Modulation}
\begin{itemize}
  \item \textbf{Monthly-resolution modulation cannot resolve daily
        variability.} $\phi$ is monthly-resolution only, so the pipeline's
        prediction varies smoothly within a month. The raw daily Pearson $r$
        against real MSL/RAD telemetry is weak ($r=-0.127$, $n=218$, not
        significant), but this figure is largely confounded by a
        documented SEP event (X5.4 flare, AR11429, 2012-03-07, one of the
        largest of solar cycle 24) that spikes RAD's raw dose rate
        $\sim\!30\times$ above baseline for a few days. This pipeline
        models SEP separately from chronic GCR, so comparing the GCR-only
        model against telemetry containing a live SEP spike biases the
        daily-resolution test for those days. Spearman's $\rho$ over the full
        series ($\rho=-0.007$) and Pearson $r$ with the event window excluded
        ($r=-0.001$, $n=209$) are both consistent with zero.
        \url{scripts/daily_phi_from_nm.py}
        tested whether giving the model genuine daily-resolution modulation
        (real Oulu neutron-monitor count rate, independently calibrated to
        $\phi$ on a 1995--2011 window that never touches the MSL cruise
        period) would rescue the correlation; it does not
        ($r=-0.127\to-0.119$), which rules out monthly $\phi$ resolution and
        points to the SEP confound as the real explanation.
        The multi-month trend the model targets remains directionally and
        quantitatively consistent regardless. Daily-resolution claims should
        be read as reflecting the monthly trend only.
  \item \textbf{ISS cross-check tests modulation shape only.} The ISS orbits
        inside Earth's magnetosphere, attenuating GCR flux by an amount not
        modeled here (no geomagnetic cutoff/trapped-radiation transport).
        The multi-year cross-check is restricted to trend correlation and
        does not validate the pipeline for near-Earth or
        trapped-radiation environments.
  \item \textbf{Force-field approximation.} Ignores charge-sign-dependent
        drift/diffusion anisotropy; during solar minimum can modify the GCR
        spectrum below $\sim\!500$~MeV/n by $10$--$20\%$ \citep{potgieter2013}.
        This is within the propagated LIS-normalization uncertainty but not
        explicitly modeled.
  \item \textbf{Missing species.} Only H, He, C, O, Si, Fe are modeled.
        Magnesium's contribution is estimated at $<3\%$ of total dose (scaled
        from Boschini 2020 Mg LIS normalization); all unmodeled species
        (Mg, N, Ne, S, Ca, ...) combined are estimated at $<5\%$.
  \item \textbf{Anomalous cosmic rays.} Not modeled; can contribute below
        $\sim\!100$~MeV/n during solar minimum.
\end{itemize}

\subsection{Biological and Risk Model}
\begin{itemize}
  \item \textbf{NASA 2023 risk framework.} This pipeline implements
        Cucinotta (2013); NASA's 2023 update introduces age/sex-dependent
        career limits and revised tissue weighting factors not yet
        implemented here.
  \item \textbf{High-LET extrapolation uncertainty.} EAR/ERR were fit
        primarily to low-LET atomic-bomb-survivor data with high-LET scaling
        from rodent experiments; no direct human epidemiological data exist
        for the HZE component of GCR. DDREF, bridging low-dose-rate space
        exposure and the acute exposures in the underlying datasets, is
        among the most debated parameters in space radiation risk assessment
        \citep{durant2013}.
  \item \textbf{Non-cancer risks.} Only cancer mortality is modeled; central
        nervous system (CNS)
        effects, cataracts, cardiovascular risk, and heritable effects are
        not quantified (CNS effects have drawn attention from rodent studies
        showing cognitive impairment below the cancer threshold
        \citep{cucinotta2014cns}).
  \item \textbf{Sex and age scope.} Results shown for 35-year-old male/female;
        a 25-year-old female has $\sim\!1.5$--$2\times$ higher REID than a
        45-year-old male for the same exposure (more remaining life-years,
        high breast/ovarian tissue weighting). The pipeline supports
        arbitrary age/sex; broader demographic analysis is future work.
\end{itemize}

\section{Supplementary Tables}

\begin{table}[t]
  \centering
  \caption{Normalized Spearman $\rho^2$ sensitivity indices ($N=500$,
           35-year-old male, 259-day transit, $16\gcm$ Al). Values sum to 1
           within each column by construction. Radiobiological parameters
           dominate REID variance; physics parameters dominate $D$ and $H$
           variance.}
  \label{tab:sensitivity}
  \small
  \begin{tabular*}{\tblwidth}{@{}LCCC@{}}
    \toprule
    Parameter & REID & $D$ & $H$ \\
    \midrule
    \code{phi\_scale}       & 0.080 & 0.793 & 0.080 \\
    \code{LIS\_norm}        & 0.001 & 0.188 & 0.002 \\
    \code{cross\_sec}       & 0.020 & $<$0.001 & 0.008 \\
    \code{neutron\_H}       & 0.009 & 0.001 & 0.007 \\
    \code{hze\_norm\_scale} & 0.008 & 0.014 & 0.009 \\
    \code{Q\_factor}        & 0.749 & 0.001 & 0.886 \\
    \code{DDREF}            & 0.051 & 0.002 & 0.003 \\
    \code{ERR\_scale}       & 0.040 & 0.001 & 0.001 \\
    \code{EAR\_scale}       & 0.043 & $<$0.001 & 0.005 \\
    \bottomrule
  \end{tabular*}
\end{table}

\begin{table}[t]
  \centering
  \caption{Value-of-information sweep: reduction in p95 REID from halving
           one parameter's uncertainty at a time (baseline p95 $=20.0\%$).}
  \label{tab:vov}
  \small
  \begin{tabular*}{\tblwidth}{@{}LCCC@{}}
    \toprule
    Parameter & Narrowed p95 & $\Delta$ p95 & \% of baseline p95 \\
    \midrule
    \code{Q\_factor}       & 13.3\% & $-$6.7 pts & $-$33.5\% \\
    \code{DDREF}           & 19.1\% & $-$0.9 pts & $-$4.5\%  \\
    \code{phi\_scale}      & 19.5\% & $-$0.5 pts & $-$2.7\%  \\
    \code{ERR\_scale}      & 19.8\% & $-$0.3 pts & $-$1.3\%  \\
    \code{LIS\_norm}       & 19.8\% & $-$0.2 pts & $-$1.1\%  \\
    \code{hze\_norm\_scale}& 19.9\% & $-$0.2 pts & $-$0.8\%  \\
    \code{neutron\_H}      & 20.0\% & $-$0.1 pts & $-$0.3\%  \\
    \code{EAR\_scale}      & 20.0\% & $-$0.0 pts & $-$0.1\%  \\
    \code{cross\_sec}      & 20.0\% & $0.0$ pts  & $0.0\%$   \\
    \bottomrule
  \end{tabular*}
\end{table}

The tool-comparison table (Ours vs.\ HZETRN, OLTARIS, Geant4, TOPAS) is
reproduced as Table~2 in the main text (\S{}Comparison with Existing Tools
and Reproducibility) rather than duplicated here.

\bibliographystyle{cas-model2-names}
\bibliography{references}

\end{document}
